\subsection{Training and Validation Architecture}

The simulation framework was operationalized through a strict three-stage computational pipeline: initialization, multi-core vectorized training, and hyperparameter stabilization. By serving as the training environment, the Agent-Based Model (ABM) allowed the Deep Reinforcement Learning (DRL) agent to thoroughly explore the massive 21-dimensional policy space through high-volume trial and error—a process that would have been computationally prohibitive or unethical in the real world \citep{TianReview2024}.

\paragraph{Phase 1: Empirical Initialization and Baseline Calibration}
First, the model was initialized using baseline data gathered from key-informant interviews with officials from all seven barangays. The environment was explicitly loaded with specific demographic profiles derived from municipal audits \citep{Villanueva2021, Yazawa2025}. For instance, \textit{Poblacion} was initialized with 1,534 households and a baseline compliance of 2\%, whereas \textit{Binuni} started with 507 households and a 15\% compliance rate. 

This step was critical to ensure the ABM accurately captured the real-world heterogeneity facing the Local Government Unit (LGU). The unobservable behavioral parameters (e.g., household sensitivity to fines) were systematically tuned until the model's unguided baseline accurately reproduced the current-day municipal compliance rate of 10–15\% \citep{Jimenez2025}. Validating this "Status Quo" baseline ensures the ABM is a credible socio-environmental system before proceeding to AI policy testing \citep{Ma2023}.

\paragraph{Phase 2: Multi-Core Vectorized Environments}
Training a DRL agent within an ABM requires a massive computational pipeline. Processing thousands of interacting household agents over multi-year horizons creates a severe computational bottleneck. To overcome this, the training architecture eschewed standard sequential loops (such as \texttt{DummyVecEnv}) in favor of true multiprocessing using a \texttt{SubprocVecEnv} wrapper.

The \texttt{Stable-Baselines3} training script spawns four parallel CPU processes ($N_{cpu} = 4$), with each core running its own isolated instance of the \texttt{BacolodGymEnv}. This architectural decision drastically improves sample efficiency by allowing the PPO agent to collect diverse environmental trajectories simultaneously, stabilizing gradient updates and significantly reducing the wall-clock convergence time of the 50,000-timestep training horizon \citep{Dey2025}.

\paragraph{Phase 3: Hyperparameter Optimization and Trajectory Rollouts}
Given the high-dimensional nature of the 21-lever continuous action space, the Proximal Policy Optimization (PPO) model was calibrated with highly specific hyperparameters to prevent policy collapse. 

To ensure the neural network gathered sufficient environmental feedback before updating its weights, the rollout buffer was expanded (\texttt{n\_steps = 1024}) coupled with a larger batch size (\texttt{batch\_size = 128}). This guarantees that the agent observes a full quarterly trajectory before executing a gradient update. Furthermore, an Entropy Coefficient (\texttt{ent\_coef = 0.01}) was injected into the loss function. This mathematical constraint actively discourages premature convergence, forcing the AI to maintain a healthy rate of exploration before "locking in" to a finalized budget allocation strategy \citep{Raffin2021}.

\vspace{0.5cm} 
\begin{algorithm}[H]
\caption{The Multi-Core DRL Training Pipeline}
\label{alg:drl_training_english}
\begin{algorithmic}[1]
\State \textbf{Input Constraints:} 
\State \parbox[t]{0.9\linewidth}{$\rightarrow$ The \texttt{BacolodGymEnv} simulation, $N_{cpu} = 4$ processing cores, and a target horizon of 50,000 algorithmic timesteps.\strut}
\Statex
\State \textbf{Step 1: Environment Vectorization (Multiprocessing)}
\State \parbox[t]{0.9\linewidth}{$\rightarrow$ \textit{Sub-Process Spawning:} Fork the main Python thread into 4 independent worker processes.\strut}
\State \parbox[t]{0.9\linewidth}{$\rightarrow$ \textit{Instantiation:} Load a complete, identical instance of the municipality (thousands of households, budgets, and variables) into each of the 4 processing cores.\strut}
\State \parbox[t]{0.9\linewidth}{$\rightarrow$ \textit{Monitoring:} Wrap the vectorized environment in a \texttt{VecMonitor} to systematically log episode rewards and simulation lengths to Tensorboard.\strut}
\Statex
\State \textbf{Step 2: Neural Network Initialization (PPO)}
\State \parbox[t]{0.9\linewidth}{$\rightarrow$ Initialize the Actor-Critic Multi-Layer Perceptron (\texttt{"MlpPolicy"}).\strut}
\State \parbox[t]{0.9\linewidth}{$\rightarrow$ Apply the stabilized hyperparameters: Learning Rate ($0.0003$), Trajectory Rollout ($1024$ steps), and Entropy Coefficient ($0.01$) to ensure persistent exploration.\strut}
\Statex
\State \textbf{Step 3: Execution of the Learning Loop}
\State \parbox[t]{0.9\linewidth}{$\rightarrow$ \textit{Parallel Data Collection:} The Mayor Agent simultaneously plays out different budgetary strategies across all 4 municipal simulations.\strut}
\State \parbox[t]{0.9\linewidth}{$\rightarrow$ \textit{Batch Updating:} After every 1,024 steps, aggregate the successes and failures from all 4 cores. Calculate the neural gradients and update the Mayor's underlying policy weights.\strut}
\State \parbox[t]{0.9\linewidth}{$\rightarrow$ Repeat this cycle continuously until the 50,000 timestep threshold is reached.\strut}
\Statex
\State \textbf{Step 4: Serialization}
\State \parbox[t]{0.9\linewidth}{$\rightarrow$ Serialize and export the fully trained neural network weights (\texttt{hudrl\_bacolod.zip}) to disk for deterministic evaluation and real-time visualization.\strut}
\end{algorithmic}
\end{algorithm}

\subsection{Hyperparameter Configuration and Network Architecture}

To effectively navigate the complex, 21-dimensional continuous action space of the municipal budgeting environment, the Proximal Policy Optimization (PPO) algorithm—operating within the system's defined heuristic guardrails—was carefully calibrated. The hyperparameters, outlined in Table \ref{tab:hyperparameters}, were selected to balance sample efficiency, gradient stability, and persistent exploration throughout the training lifecycle \citep{Raffin2021}.

The underlying policy and value functions were approximated using a Multi-Layer Perceptron (\texttt{MlpPolicy}) architecture, specifically configured with two hidden layers of 64 units each. This relatively lightweight architecture prevents overfitting while maintaining sufficient representational capacity to map the high-dimensional municipal state space to specific budgetary actions. To overcome the computational bottleneck of simulating thousands of household agents, the environment utilized a \texttt{SubprocVecEnv} wrapper distributed across four CPU cores. Training was conducted over a total horizon of 500,000 timesteps, providing the agent with extensive environmental interactions to refine its resource allocation strategy \citep{Mousavi2021}.

To optimize this architecture, the learning loop was governed by a conservative learning rate of $0.0003$, which, when paired with PPO's default clipping range of $0.20$, prevents destructive, excessively large policy updates \citep{Raffin2021}. During each iteration, the multi-core environment collects trajectories across a horizon (\texttt{n\_steps}) of 1024 timesteps per core, yielding a comprehensive batch of 4,096 experiences per update cycle. This data is then processed in smaller mini-batches of 128 over 10 optimization epochs. Furthermore, the Generalized Advantage Estimation (GAE) parameter ($\lambda = 0.95$) was utilized to strike an optimal balance between bias and variance when estimating the advantage function, ensuring that the credit assignment for specific budgetary interventions remains accurate across lengthy policy rollouts.

Beyond immediate gradient updates, the agent's strategic foresight is controlled by a discount factor ($\gamma$) of $0.99$, compelling the model to prioritize long-term municipal compliance and systemic sustainability over immediate, short-term reward spikes. Because the continuous action space is vast, premature convergence to sub-optimal local minima is a significant risk. To mitigate this, an entropy coefficient (\texttt{ent\_coef}) of $0.01$ was integrated into the loss function. This introduces a mathematical penalty for overly deterministic behavior, effectively forcing the agent to maintain a healthy degree of exploration before locking into a finalized policy structure \citep{Abdallah2021, Raffin2021}.

Finally, a critical element driving this exploration is the multi-objective reward function, scaled by four distinct weights ($\omega_1 = 1.0, \omega_2 = 5.0, \omega_3 = 2.0, \omega_4 = 0.5$). These predefined scalar values act as strict pedagogical boundaries, explicitly prioritizing critical municipal outcomes—such as heavily weighting $\omega_2$ to strongly incentivize regulatory compliance and budget efficiency—while dynamically suppressing less critical infractions. To ensure that the resulting learned policy is statistically robust and not merely an artifact of favorable initialization, the entire 500,000-timestep training pipeline was rigorously evaluated across five distinct random seeds (Seeds 0--4).

\begin{table}[htbp]
\centering
\caption{PPO Hyperparameters and Training Configuration}
\label{tab:hyperparameters}
\begin{tabular}{l l}
\toprule
\textbf{Parameter} & \textbf{Value} \\
\midrule
Network Architecture & \texttt{MlpPolicy} (2 layers, 64 units) \\
Learning Rate & $0.0003$ \\
Batch Size & $128$ \\
Epochs (\texttt{n\_epochs}) & $10$ \\
Horizon (\texttt{n\_steps}) & $1024$ \\
Discount Factor ($\gamma$) & $0.99$ \\
Entropy Coefficient (\texttt{ent\_coef}) & $0.01$ \\
Clip Range & $0.20$ (Default) \\
GAE Parameter ($\lambda$) & $0.95$ (Default) \\
Reward Weights ($\omega_1, \omega_2, \omega_3, \omega_4$) & $1.0, 5.0, 2.0, 0.5$ \\
Total Timesteps & $500,000$ \\
Compute Infrastructure & 4 CPU Cores (\texttt{SubprocVecEnv}) \\
Random Seeds Evaluated & $5$ (Seeds 0--4) \\
\bottomrule
\end{tabular}
\end{table}